% NAF 5176, batch 2 — Gallica views 21–40.
% Pass: Opus 5.5 (claude-opus-5-5), 2026-09-22 — first pass, unchecked against the views by a human.
%
% The views. Greyscale scans, landscape (about 8730 × 6360–6410): each view is
% the volume lying open. Many leaves are mounted on guards, several carry
% pasted slips (béquets) that Gallica photographs twice, once folded down and
% once lifted. Text is given once, on the view where it reads whole; the other
% exposure gets a \page{} with a note only.
%
% What the twenty views hold. No piece is signed or dated on the leaf; nothing
% below is an attribution.
%
%   v. 21–29  French optical propositions in one fast cursive with many
%             deletions: « I Proposition » (given the refraction of one ray,
%             the refraction of all others on the same surface; a Corollaire),
%             « II Proposition » (measuring the refraction of a crystal prism
%             ABC with two pinnules), Corollaire citing « la 49 prop. du 1
%             livre de mes Coniques » and « la 51 » (the foci of the
%             hyperbola), « Proposition III » (hyperbola: sines in the ratio
%             of the transverse axis to the focal distance), « Proposition
%             IV » (the same for the ellipse), a Corollaire, and « V
%             Proposition » (the surface that makes rays parallel, convergent
%             or divergent: hyperbolic conoids and ellipsoidal spheroids for
%             lenses « pour les lunettes de longue veue »). v. 21 is the same
%             leaf as v. 20 re-shot with its two slips lifted; v. 22–24 and
%             25–26 are pairs of exposures of the same openings. The slips
%             carry the figures, drawn with serif capitals so that they look
%             printed. The proposition V breaks off at the foot of v. 26 on
%             the catchword « Je dis » and resumes at the head of v. 29 (« Je
%             dis en 2 lieu »): between them, v. 27 is a torn strip in the
%             same hand, a devotional fragment on light reflected back to its
%             source (« Luminaire … soleil divin »), and v. 28 is the blank
%             back of that strip. v. 29 ends mid-page.
%   v. 30     Blank opening (the back of v. 29, the recto showing through; a
%             blank leaf numbered 22). Skipped.
%   v. 31–32  In the same hand, smaller: the tangent to an ellipse by
%             « adæqualité » in the species notation (B in G, Aq, « bis »),
%             compared with Apollonius's construction; the writer's « ma
%             methode », its « 1er » and « 2 cas », the centres of gravity
%             « dont j'ay envoyé les exemples à Mr Roberval », and a question
%             proposed « à soudre ». v. 32: « Mr Fermat dans sa derniere let.
%             de Juillet 1638 » (the only date read in the batch), « Mr de
%             Cartes », « page 344. l. 3. », a curve of degree ten in x and y,
%             a figure of six semicircles, and four questions in Latin
%             numbered 1–4 (six semicircles cut by a perpendicular, sum given
%             / maximum or minimum; the surface of a scalene cone; right
%             triangles in numbers whose area is a square, or a proof of
%             impossibility). Much of v. 32 is fast and only partly read.
%   v. 33–35  A letter or discourse in French, same hand, heavily corrected:
%             the speed of light compared with a horse, birds, an arrow, a
%             musket ball, the earth's rotation (7200 lieues round), the
%             stars; the gifts of the glorified body (agilité, clarté,
%             subtilité, impassibilité); light as the image of its source;
%             the dip of the horizon from an eye 6 ft (then « une toise »)
%             high, giving the earth's radius 1250 lieues and circumference
%             7857 1/7 lieues (lieue = 2500 toises, « depuis la Bastille
%             jusques au Chasteau du Bois de Vincennes », doubtful). It ends
%             on v. 35 L: « … que j'aye employé quelque temps à la rendre
%             françoise pour ceux qui n'entendent point d'autre langue »,
%             i.e. a preface to an optics put into French. No signature, no
%             addressee named. A long vertical marginal addition on v. 33 is
%             almost wholly illegible.
%   v. 36     « Propositions » (same hand, more composed): rays falling on
%             a plano-convex crystal CFDGC, all meeting at I; then
%             « Proposition » « Expliquer pourquoy se rompent les rayons ».
%   v. 37     Blank opening (back of v. 36; a blank leaf numbered 29).
%             Skipped.
%   v. 38     « Proposition XXII. Expliquer pourquoy si l'air fait faire
%             refraction à la lumiere, & si elle peut passer par le vuide »,
%             stopping mid-page on « Quant à l'autre partie, ».
%   v. 39     A modern separator leaf: « Pascal / autographe / et copies »,
%             « Libri 1848 », « FR. nouv. acq. 5176 ». Given as a note only.
%   v. 40     A new hand (round, regular, sloping; brown, creased paper whose
%             top right corner hides line-ends): « Regle auxquelles se
%             peuvent rapporter les partis », the division of stakes in an
%             interrupted game (9 or 13 → 11; « Exemples suivant ces Regles »,
%             « En 2 Jeux combien vaut le 1er », 6 escus; two numbered Regles;
%             « Exemple — on joue en 3 jeux 8 escus contre 8 autres » with
%             shares in sixteenths). Two lines in a heavier ink look like a
%             later addition; a pasted marginal slip of about ten lines was
%             not read. It runs to the foot of the page and continues, if at
%             all, on v. 41 (batch 3).
%
% Numbers on the leaves. One bold, round ink series stands at the top right of
% every recto AND of every pasted slip, and it is continuous across the batch
% once the slips are counted as leaves: 15 (v. 21, underlined; batch 1 reads
% 13 and 14 on the two slips of this leaf at v. 20), 16 and 17 (the two slips
% of v. 22), 18 (v. 24), 19 (the slip of v. 25), 20 (the strip of v. 27), 21
% (v. 29), 22 (v. 30, blank), 23 (v. 31), 24 (v. 32, blank half of a
% bifolium), 25 (v. 33, underlined), 26 (v. 34), 28 (v. 36), 29 (v. 37,
% blank), 30 (v. 38), 31 (v. 39, the separator leaf), 32 (v. 40). 27 was not
% seen: it should fall on the slip of v. 35 R, whose numbered corner is not
% visible in either exposure. A thin « 20 » above the torn top edge of the
% leaves at v. 21 and v. 24 stands on a paper behind them, apparently the
% strip numbered 20. This series behaves exactly like a library foliation
% (one number per leaf and per slip, 15→32 here, reaching 32 at view 40 in a
% volume catalogued at 42 ff.). Batches 1 and 3 read the same series (1–15,
% 32–42), and the 1888 certificate's list of blank leaves (« Les Feuillet 22,
% 24, 29, 35, 36 sont blancs », v. 7) falls on blank leaves of this series at
% all five points. It is nevertheless penned, not pencilled, so, as for Latin
% 17859 and Français 9115, no \folio{} is written in any of the three batches
% until a human has looked at the leaves; the mapping, if accepted, is
% v. 21 = f. 15r, v. 24 = 18r, v. 27 = 20r, v. 29 = 21r, v. 31 = 23r, v. 33 =
% 25r, v. 34 R = 26r, v. 36 R = 28r, v. 38 = 30r, v. 40 = 32r (slips 16, 17,
% 19 on v. 22 and v. 25). A second, finer and older series, irregular, stands
% beside it: 16 (v. 21), 18 (v. 24), 17 (v. 27), 18 (v. 29), 21 (v. 31), 12
% (v. 33), 13 (v. 34), 40 above 28 (v. 36), 39 above 30 (v. 38), 23 beside 32
% (v. 40). On v. 39 a large « 35 » stands above the 31. Round stamps
% « BIBLIOTHÈQUE NATIONALE MANUSCRITS » / « R.F. » on most written pages and
% on the slips.
%
% The hand of v. 21–38. One fast cursive (v. 31–32 smaller, v. 36 more
% composed): the ampersand is a large looped C, rendered \&; « en » is
% written with a sign like a 3 and set « en »; long s set s throughout; a
% long horizontal is often a genuine deletion here. Figures are not drawn:
% their lettering is given in notes. Names of points on v. 31–32 carry a bar,
% set \overline{}; the Greek-looking letters there are read ζ, δ, ε from the
% figure. The equality sign of v. 32 (a reversed & under the first member)
% is set \propto, an inference noted there. Spelling of the period kept
% throughout (mesme, estoille, sçavoir, adiouster, soudre, lors que, c'est
% àdire, parceque, pourceque).
%
% Continuities with the neighbouring batches. v. 21 is v. 20 re-shot: the
% text of the leaf of v. 20 is given here whole. v. 40 runs on to its foot;
% v. 41 may continue it.

\documentclass[11pt,a4paper]{article}
% The preamble shared by every transcription of the mathematicians' manuscripts in Gallica.
%
% It rests on one idea: **uncertainty compounds, it does not resolve.** A
% transcription that smooths over a doubtful word has destroyed the only thing
% it was made for — a reader must be able to tell, at every point, what was on
% the page and what was supplied. The apparatus macros below are therefore the
% critical apparatus, not decoration.
%
% The same commands are recognised by `scripts/render.mjs`, which produces the
% site's reading view, and by `scripts/tei.mjs`. A macro added here must be
% added there too, or rendering fails loudly — which is the wanted behaviour.
%
% Comments here are English, like the rest of the repository. The documents
% this preamble sets are French, and that is deliberate: see the skills.

\ProvidesPackage{germain}[2026/09/15 Transcriptions of mathematicians' manuscripts in Gallica]

\usepackage{amsmath,amssymb,amsthm}
\usepackage{mathrsfs}
% Kept from the parent preamble so the renderer's diagram support stays
% available; these manuscripts draw no commutative diagrams, and nothing here uses it.
\usepackage{tikz-cd}
\usepackage[english,french]{babel}
\usepackage{fontspec}

% --- Greek ------------------------------------------------------------------
% The text face stays fontspec's default, Latin Modern, which has no Greek:
% XeTeX drops every glyph it lacks with a line in the log and nothing on the
% page, and the Greek words the manuscripts quote vanished from the PDFs. So
% the Greek and Greek Extended blocks get a XeTeX character class of their own,
% and entering or leaving a run of it switches the family to CMU Serif — the
% same Computer Modern design, with polytonic Greek — and back. Only the
% family changes: a Greek word in \textit or \textbf keeps its shape and series.
% The transitions to and from every other class carry the tokens that class
% has towards an ordinary letter, so French spacing before « ; » or inside
% « » still holds after a Greek word. No grouping, so no brace or \endgroup
% can come out unbalanced; maths is left alone, since XeTeX inserts nothing
% there. Kept identical in `exercices.sty`.
\makeatletter
\newfontfamily\germainGreekfont{cmun}[
  Extension=.otf, NFSSFamily=germaingreek,
  UprightFont=*rm, ItalicFont=*ti, SlantedFont=*sl,
  BoldFont=*bx, BoldItalicFont=*bi, BoldSlantedFont=*bl]
\newXeTeXintercharclass\germain@greek
\def\germain@greekrange#1#2{%
  \@tempcnta=#1\relax
  \loop
    \XeTeXcharclass\@tempcnta=\germain@greek
    \ifnum\@tempcnta<#2\advance\@tempcnta\@ne\repeat}
\germain@greekrange{"0370}{"03FF}
\germain@greekrange{"1F00}{"1FFF}
\def\germain@greekprev{\rmdefault}
\def\germain@greekin{%
  \edef\germain@greekprev{\f@family}\fontfamily{germaingreek}\selectfont}
\def\germain@greekout{\fontfamily{\germain@greekprev}\selectfont}
\def\germain@greekjoin{%
  \XeTeXinterchartokenstate=\@ne
  \@tempcnta=\z@
  \loop
    \ifnum\@tempcnta=\germain@greek\else
      \XeTeXinterchartoks\germain@greek\@tempcnta=\expandafter{%
        \expandafter\germain@greekout\the\XeTeXinterchartoks\z@\@tempcnta}%
      \XeTeXinterchartoks\@tempcnta\germain@greek=\expandafter{%
        \the\XeTeXinterchartoks\@tempcnta\z@\germain@greekin}%
    \fi
    \ifnum\@tempcnta<4095\advance\@tempcnta\@ne\repeat}
% babel-french sets its own classes, and saves and restores the interchar
% state, each time a language is selected: join again after it does.
\addto\extrasfrench{\germain@greekjoin}
\addto\extrasenglish{\germain@greekjoin}
\AtBeginDocument{\germain@greekjoin}
\makeatother

\usepackage{geometry}
\geometry{a4paper,margin=25mm,marginparwidth=28mm}
\usepackage{marginnote}
\usepackage{xcolor}
\usepackage[normalem]{ulem}
\setlength{\emergencystretch}{3em}
\usepackage{hyperref}
\hypersetup{colorlinks=true,linkcolor=black,urlcolor=black,pdfborder={0 0 0}}

% --- Batch metadata ---------------------------------------------------------
% Taken as they stand by the rendering script, which shows them at the head of
% the reading view. They fix what the file claims to cover. The macro names are
% the parent project's, so that the scripts stay shared; \folder here means the
% volume's slug — `fr-9115` — and \pages counts Gallica views.

\newcommand{\folder}[1]{\gdef\grothFolder{#1}}
\newcommand{\batch}[1]{\gdef\grothBatch{#1}}
\newcommand{\pages}[2]{\gdef\grothFirst{#1}\gdef\grothLast{#2}}
\newcommand{\foldertitle}[1]{\gdef\grothFoldertitle{#1}}

% The holder's own shelfmark, printed as they write it: « Français 9115 ».
\newcommand{\shelfmark}[1]{\gdef\grothShelfmark{#1}}

% The dating, copied verbatim from the holder's catalogue — never inferred
% here. « XIXe siècle » is the BnF's; a pass that dates a leaf from its content
% says so in a \note{}, as its own inference.
\newcommand{\dating}[1]{\gdef\grothDating{#1}}

% The status notice, printed in the title block and repeated as a diagonal
% watermark on every page of the PDF. The manuscripts are in the public
% domain; what this line declares is not a right but a quality — an
% unverified first pass by a machine. The skills write the line; a file
% without it gets no watermark, so leaving it out is visible in the output.
\newcommand{\watermark}[1]{\gdef\grothWatermark{#1}}

\gdef\grothFolder{?}\gdef\grothBatch{}\gdef\grothFirst{?}\gdef\grothLast{?}%
\gdef\grothFoldertitle{}\gdef\grothShelfmark{}\gdef\grothDating{}\gdef\grothWatermark{}

% --- The résumé -------------------------------------------------------------
\newenvironment{resume}
  {\small\setlength{\parskip}{0.45em}}
  {\par\medskip\hrule\medskip}

% --- Keywords ---------------------------------------------------------------
% In English, closing the modernised reading's résumé: the modern vocabulary
% under which to search for what the leaves construct. The manifest extracts
% them to tag the volume — this line is the single source of the tags.
\newcommand{\keywords}[1]{\par\smallskip\noindent
  {\footnotesize\textsc{Keywords} --- \textit{#1}}\par}

% --- The critical apparatus -------------------------------------------------

% The Gallica view number, in the margin. This is the anchor that lets a
% disagreement over a formula be settled by looking at *one* image rather than
% a volume of 750. The site uses it to turn the facsimile as the transcription
% is scrolled. A view is one scanned image: usually one side of a leaf.
\newcommand{\page}[1]{%
  \marginnote{\footnotesize\color{gray!70}\textsf{v.\,#1}}%
  \label{page:#1}\ignorespaces}

% The BnF's pencilled foliation, where the leaf carries it and a pass can read
% it: « 198r ». This is what the literature cites — « ff. 198r–208v » — and
% the one line that turns a citation into a view. Recorded, never inferred:
% a folio a pass cannot read is not written.
\newcommand{\folio}[1]{%
  \marginnote{\footnotesize\color{gray!50}\textsf{f.\,#1}}\ignorespaces}

% The modernised reading's coarser counterpart to \page{}: which views a
% section is drawn from.
\newcommand{\pagerange}[2]{%
  \marginnote{\footnotesize\color{gray!70}\textsf{v.\,#1--#2}}%
  \ignorespaces}

% Illegible: never guessed.
\newcommand{\ill}{\textcolor{red!60!black}{[\,\ldots\,]}}

% A doubtful reading: the word is given, and flagged as uncertain.
\newcommand{\uncertain}[1]{\uline{#1}}

% An editorial addition — a missing letter, an implied word.
\newcommand{\add}[1]{\textcolor{blue!50!black}{[#1]}}

% Struck out by the writer. What was crossed out is part of the document.
\newcommand{\struck}[1]{\sout{#1}}

% The transcriber's note, on the state of the page or the mathematics.
\newcommand{\note}[1]{\par\smallskip{\small\color{gray!60!black}\textit{[#1]}}\par\smallskip}

% A marginal note of hers, distinct from the body of the text.
\newcommand{\marginal}[1]{\par\smallskip\noindent\hspace{1em}%
  {\small\color{gray!60!black}#1}\par\smallskip}

% --- Title ------------------------------------------------------------------
% The title block is French, like the documents themselves.

\AddToHook{shipout/background}{%
  \ifx\grothWatermark\empty\else
    \begin{tikzpicture}[remember picture, overlay]
      \node[rotate=52, scale=3.4, align=center, text=gray!18,
            font=\sffamily\bfseries]
        at (current page.center) {\grothWatermark};
    \end{tikzpicture}%
  \fi}

\AtBeginDocument{%
  \begin{center}
    {\large\bfseries Manuscrits de mathématiciens (Gallica) ---
      \ifx\grothShelfmark\empty\grothFolder\else\grothShelfmark\fi}\\[2pt]
    {\itshape\grothFoldertitle}\\[4pt]
    {\small vues \grothFirst--\grothLast
      \ifx\grothBatch\empty\else\ (lot \grothBatch)\fi}\\[2pt]
    \ifx\grothDating\empty\else
      {\small\color{gray!60!black}Datation du catalogue : \grothDating}\\[2pt]
    \fi
    \ifx\grothWatermark\empty\else
      {\small\bfseries\color{red!45!gray}\grothWatermark}\\[2pt]
    \fi
    {\footnotesize\color{gray!60!black}Transcription établie d'après les images IIIF de Gallica.
     Source gallica.bnf.fr / Bibliothèque nationale de France.\\
     Les lectures incertaines sont soulignées, les illisibles notées [\,\ldots\,],
     les ajouts entre crochets ; \textsf{v.} numérote les vues de Gallica,
     \textsf{f.} la foliotation au crayon de la BnF.}
  \end{center}
  \vspace{1em}}

\endinput


\folder{naf-5176}
\shelfmark{NAF 5176}
\batch{2}
\pages{21}{40}
\dating{1601-1700}
\watermark{Lecture automatique\\non vérifiée}
\foldertitle{Papiers du Père MERSENNE et de Blaise PASCAL.}

\begin{document}

\page{21}
\note{Même feuillet que la vue 20, photographié de nouveau avec les deux béquets relevés vers la gauche : le texte que les béquets cachaient à la vue 20 se lit ici en entier. En haut à droite, trois nombres : « 20 », d'un trait fin, au-dessus du bord déchiré du feuillet, sur un papier qui dépasse derrière lui (sans doute la bande numérotée 20 de la vue 27) ; sur le feuillet, « 15 » à l'encre, rond et souligné (le 5 a la forme d'un S), et à sa droite « 16 », d'un trait fin. Le « 15 » ouvre la série continue décrite dans l'en-tête ; il n'est pas donné ici comme folio, voir l'en-tête.}

I Proposition
\note{titre centré. Le texte de l'énoncé est tenu dans une grande accolade ouvrante à gauche, comme tous les énoncés de ces feuillets.}

L'inclination \& la refraction de quelconque rayon sur \struck{quelconque} \add{la} surface \struck{d'un corps diaphane} de quelconque corps diaphane qui le rompe, estant donnée ; donner la refraction de tous \uncertain{les autres} rayons inclinez sur la mesme surface.
\note{« la » est écrit au-dessus de « quelconque » biffé. L'esperluette a partout la forme d'un grand C bouclé ; elle est rendue « \& ».}

\struck{Soit} \uncertain{Et que} le plan, où est le rayon \struck{AEB} ADB soit la commune section de quelconque corps \& d'\ill{} tel plan qui le rompra, \uncertain{\& partant} que la droite AEB soit la commune section de la surface. Et que la refraction du rayon incliné FE, se \uncertain{face} en EG. Ayant tiré la perpendiculaire CED, CEF soit l'angle de l'inclination du rayon FE, \& KEG soit l'angle de la refraction de \uncertain{ce mesme} rayon. \struck{\ill{}} Si soit choisi tel autre rayon qu'on voudra \struck{\ill{}} comme HE incliné suivant l'angle CEH, lequel aura sa refraction \ill{} \uncertain{autre que la} premiere.
\note{un mot bref est écrit sous la ligne, sous « rompra », sans signe d'insertion ; il n'a pas été lu. Le début de la phrase (« Et que le plan, où est le rayon ADB soit la commune section… ») est donné tel qu'il se lit ; la construction en est incertaine.}

Soit \add{descrit} du centre E \struck{\ill{}} \& de quelconque intervalle EB \struck{descrit} le demi cercle ACB, dont la circonference coupe le rayon FE au point F, \& le rayon HE au point H ; \& après avoir mené \uncertain{par} \ill{} tiré la ligne FI parallele à la ligne AB, soit tirée IG parallele à CE, laquelle rencontre le rayon rompu EG au point G.
\note{« descrit » est écrit au-dessus de la ligne, avec un signe d'insertion, et le même mot est biffé plus loin.}

\struck{Et du mesme centre E, de l'intervalle EG soit descritte la circonference du cercle LGD.}
Et que la circonference du cercle LGD descritte du centre E \& de l'intervalle EG, rencontre les diametres \struck{prolongez} CE, BE prolongez aux points L \& D. Et après avoir tiré HM parallele \struck{\ill{}} à BA, soit descritte MN parallele à CD qui rencontre la circonference au point N. La ligne EN sera la refraction du rayon HE qui est rompu au point E, \& qui fait l'angle rompu DEN.
\note{la phrase biffée l'est d'un long trait continu. « DEN » est suivi de tirets de remplissage jusqu'au bout de la ligne.}

On trouvera de la mesme façon \struck{\ill{}} la refraction de tous les autres rayons.
\note{sous la ligne, un long trait ondulé qui la souligne jusqu'à la marge. Dans la marge de gauche, en face de cette ligne, un tiret. Le cachet rond « BIBLIOTHÈQUE NATIONALE MANUSCRITS » est apposé ici.}

Corollaire.

D'où il s'ensuit que deux diaphanes estant \uncertain{comme} AEB \& LDB, dont \ill{} plus rare \& \uncertain{l'autre} plus dense, leurs angles d'inclination \& de refraction doivent estre comparez ensemble suivant la raison \struck{\ill{}} de la droite GE à la droite EF.

II Proposition.

Estant donné quelconque diaphane trouver la refraction du rayon qui traverse \ill{} ledit diaphane.
\note{énoncé tenu dans une accolade. La fin de la première ligne et le premier mot de la seconde (« de l'air » ? « partie » ?) n'ont pas été lus.}

Soit fait le triangle rectangle X \uncertain{d'un morceau} dudit diaphane suivant le triangle ABC, dont l'angle droit est en B, \& que les 2 autres angles A \& C \struck{soient} tels qu'on voudra, pourveu que l'angle A soit plus aigu que l'angle C. Et ayant mis le costé BC sur \struck{\ill{}} \add{le} plan horizontal DBE, que le rayon du soleil, ou de tel \ill{}, soit conduit par les 2 pinnules, dont les \uncertain{trous} FG soient fort \ill{}, le rayon FG \struck{\ill{}} tombant sur le cristal AB le transperce droit de G en I, parcequ'il \struck{\ill{}} tombe à plomb, \& se rompra au point I sur la surface oblique AC, pour aller au point E par le rayon droit IE, parceque le milieu AOE est l'air qui est plus rare, \& plus subtil que le cristal ABC.
\note{la page s'arrête ici, au milieu de la démonstration ; la suite est au verso, vue 22, page de gauche.}

\note{Les deux béquets, relevés, sont collés au bord gauche du feuillet. Le béquet du haut, qui à la vue 20 recouvre le haut de la page, montre ici son dos : la figure de la proposition I y transparaît à l'envers — un cercle de centre E, un diamètre horizontal et une perpendiculaire, des rayons, des lignes pointillées, lettres A, B, C, D, E, H, I, K, M, N, O. Le béquet du bas montre ici sa face : la figure de la proposition II, d'apparence imprimée, à lettres capitales à empattements — le triangle X aux côtés hachurés, et la construction lettrée A, B, C, D, E, F, G, H, I, K, L, M, N, O, P, Q, avec un arc de cercle de F à K par D. Le cachet de la bibliothèque est répété, à l'envers, sur les deux béquets.}

\page{22}
\note{Ouverture : à gauche, le verso du feuillet de la vue 21 ; à droite, le recto du feuillet suivant. Les deux pages portent chacune un béquet rabattu sur le bas du texte. Le texte n'est donné qu'une fois, là où il se lit en entier : la page de gauche à la vue 23, où son béquet est relevé, la page de droite à la vue 24, où les deux le sont. Ici ne se voient que les deux béquets. Celui de gauche porte, à l'encre, en haut à gauche, le nombre « 16 », et une figure d'apparence imprimée, à lettres capitales à empattements : une hyperbole de sommet B, son axe vertical portant les points F, D, H, B, E, O, I, et, autour du point C de la courbe, les droites et l'arc de cercle de la construction, lettres C, G, K, L, M, N, P, Q. Celui de droite porte « 17 » en haut à droite et une figure de même facture : une ellipse d'axe BF, avec les points H, E, O, I, D sur l'axe, C sur la courbe, et la même construction, lettres G, K, L, M, N, P, Q. Le cachet « BIBLIOTHÈQUE NATIONALE R.F. » est apposé sur chacun. Les nombres « 16 » et « 17 » des béquets appartiennent à la série continue de l'en-tête. En haut à droite de la page de droite, « 20 » et « 18 » ; voir la vue 24.}

\page{23}
\note{La même ouverture que la vue 22, le béquet de gauche relevé vers la gauche : on n'en voit plus que le dos, la figure transparaissant à l'envers. La page de gauche, verso du feuillet de la vue 21, se lit ici en entier ; elle poursuit la proposition II. Aucun nombre en tête de cette page.}

\uncertain{Cela} estant fait, \& \uncertain{ayant} tiré sur la droite IC \struck{\ill{}} \add{la} perpendiculaire DIH, \& \ill{} l'angle DIF pour l'inclination du rayon FI, \& HIE soit l'angle rompu, \uncertain{de prolonger} la perpendiculaire HID, sert de costé tant à l'angle d'inclination, qu'à l'angle rompu du rayon.
\note{« la » est écrit au-dessus d'un mot biffé. Le mot qui suit le second « \& » n'a pas été lu. La construction de la phrase est donnée telle qu'elle se lit.}

Et partant, si du centre I, de l'intervalle ID est descrit la circonference DF se rencontrant avec le rayon FI au point F, \& après tiré FK parallele à IC se rencontrant ladite circonference au point K, \& puis menée KL parallele à ID, \& au rayon rompu IE au point L, lon aura la raison de la refraction de tous les rayons qui \add{viendront de l'air} \struck{tomberont} sur ce diaphane, pour ce \uncertain{qu'elle} sera tousjours la mesme que celle que \struck{\ill{}} lon aura trouvée \struck{\ill{}} de la droite LI à la droite ID ou IF.
\note{« viendront de l'air » est écrit au-dessus de « tomberont » biffé.}

Corollaire

Ayant tiré MI qui face l'angle MIC egal à l'angle CIE, parceque si des deux foyers de l'hyperbole (suivant la 49 prop. du 1 livre \uncertain{de mes} Coniques) deux lignes droites s'inclinent à un mesme point de la section, elles feront deux angles egaux avec la \struck{contin} ligne \uncertain{contingente} qui touche. Supposé que \struck{\ill{}} les 2 points M \& E soient les deux foyers de l'hyperbole, \& que le point I soit dans la section ; la droite IC touchera la section au mesme point I.
\note{« de mes Coniques » : entre « de » et « Coniques », un mot bref, à initiale bouclée, qui se lit « mes » ou « ses ». La même référence revient plus bas (« la 51. … 1 livre des Coniques », « mon 2 l. des Con. ») : l'auteur semble renvoyer à des Coniques de sa main, par livre et par proposition. Lecture à vérifier. Le chiffre 49 est net.}

Partant, si après avoir pris \add{dans IE} la droite \struck{\ill{}} IN egale à IN, lon fait EO egale à EN, \& qu'après avoir divisé MO en 2 parties egales au point P, lon face EQ egale à MP, d'autant que suivant la 51. \uncertain{de mes} 1 livre des Coniques si dans l'hyperbole deux droites s'inclinent \struck{descrivent} du foyer \struck{\ill{}} \& à un mesme point de la section la plus grande surpasse la moindre de la longueur de l'axe transversant ; \& les 2 points P \& Q seront les \uncertain{sommets} \uncertain{contreposez} de l'hyperbole, \& partant l'axe transversant sera la \struck{ligne} \add{droite} PQ ; \struck{\ill{}}
\note{« dans IE » est écrit au-dessus de la ligne, en face d'un mot noirci par une rature. « IN egale à IN » : les deux fois la leçon est IN, la seconde peut-être pour IM ; transcrit tel qu'écrit. « droite » est écrit au-dessus de « ligne » biffé.}

Or tout cela estant fait, il sera aysé de descrire \struck{\ill{}} l'hyperbole, suivant la 20, 26 \& les autres \uncertain{de mon} 2 l. des Con.
\note{« 2 l. des Con. » : abréviation de « 2 livre des Coniques ».}

Proposition III.

Dans toute hyperbole, si de quelconque point pris dans la section, l'on tire 2 droites, l'une par le foyer exterieur, l'autre parallele à l'axe \& que par le mesme point on mene une troisiesme droite perpendiculaire à la tangente de la section, le sinus de l'angle fait par la parallele à l'axe \& par la perpendiculaire sera au sinus de l'angle fait par la droite tirée par le foyer, \& par la mesme perpendiculaire, comme l'axe transversant de la section à la distance d'entre les foyers.
\note{énoncé tenu dans une accolade.}

Soit telle hyperbole qu'on voudra ABC, dont le sommet est B, le foyer exterieur D, le centre E, l'axe transversant BF. De quelconque point, comme C, soient menées les droites CD par le foyer D, \& CG parallele à BD ; \& après avoir tiré la tangente CH, que CI luy soit perpendiculaire.
\note{la page s'arrête ici ; la démonstration se poursuit en tête de la page de droite, vue 24.}

\page{24}
\note{La même ouverture, les deux béquets relevés ; celui de droite, rabattu vers la gauche, couvre à présent le bas de la page de gauche. La page de droite se lit ici en entier ; elle poursuit la démonstration de la proposition III. En haut à droite, « 20 », d'un trait fin, tout en haut, sans doute le nombre d'un autre papier qui dépasse comme à la vue 21, et plus bas, sur le feuillet, « 18 », rond. Le cachet de la bibliothèque est apposé dans la marge.}

\struck{Et puis} De plus, du centre C, de l'intervalle CI soit descritte la circonference du cercle qui coupe la droite CG au point G, \& la droite DC au point K. Ayant \struck{\ill{}} \add{faict} les droittes GL, KM perpendiculaires à CI, je dis que le sinus GL est au sinus KM, comme BF à ED, c'est àdire comme l'axe transversant de l'hyperbole ABC, à l'intervalle d'entre les foyers.
\note{« Et puis » est biffé au-dessus de la première ligne. « faict » est écrit au-dessus d'un mot biffé.}

Du point I soit tirée la perpendiculaire IN sur la droite DK ; \& du point C soit tirée CO perpendiculaire à l'axe ; \& après avoir descrit la tangente CH, \& l'avoir \uncertain{conjointe} à CE, que EP soit menée sur DO \uncertain{en sorte} qu'elle soit parallele à IC, \& qu'elle coupe la tangente au point Q.
\note{après « DO », une petite croix, peut-être un renvoi ; aucune note ne lui répond sur la page.}

A cause des droites egales CI, CG, \& des angles egaux CIO, \& ICG, \& de l'angle droit en O \& L, la droite CO sera egale à GL. Et semblablement, à cause des egales CI, CK \& de l'angle commun C, \& des angles droits M \& N, la droite IN sera egale à KM. \struck{Partant} Par consequent, comme CO sera à IN, ou DC à DI, ou DP à DE, ainsi GL sera à KM.

Mais pourceque CI, EP sont perpendiculaires à la tangente CH, \& que les angles DCH, HCE sont egaux, par la susdite 49 prop. du 1, \& que les 2 angles faits au point Q sont droits, CP sera egale à CE, \& DP à l'axe transversant BF. \struck{Dor}

Donc le sinus GL sera au sinus KM, comme l'axe transversant BF de l'hyperbole ABC, est à la distance \struck{des foyers} DE des foyers, ce qu'il falloit demonstrer.
\note{un mot bref est écrit au-dessus de « des foyers » biffé ; il n'a pas été lu.}

Proposition IV.

Dans toute Ellipse, si de quelconque point, qui ne soit pas au milieu de la section, l'on meine 2 droites, l'une par le foyer le plus eloigné, l'autre parallele à l'axe ; \& qu'on conduise \uncertain{encore} du mesme point une troisiesme droite \struck{qui} parallele à la tangente, les sinus des angles faits par la parallele à l'axe, \& par la droite menée par le foyer avec la perpendiculaire, seront entre eux comme \struck{l'axe} le plus grand axe transversant de la section à la distance des foyers.
\note{énoncé tenu dans une accolade. « Ellipse » est récrit sur un premier tracé. La troisième droite est dite « parallele à la tangente », puis la fin de l'énoncé parle de « la perpendiculaire » : transcrit tel qu'écrit ; la démonstration qui suit la mène perpendiculaire.}

Soit quelconque Ellipse ABC, dont \struck{l'axe transversant} le plus grand axe transversant est BF, les foyers DE : \& soit pris quelconque point C dans la section pourveu qu'il ne soit point au milieu (qui est là où se termine l'un \& l'autre axe). Soit tirée CD par le foyer le plus eloigné D, \& CG parallele à l'axe BF.

Et après avoir tiré la tangente CH, \struck{soit menée} que CI luy soit perpendiculaire \struck{\ill{}} \& rencontre l'axe en I. Du centre C, de l'intervalle CI soit descritte la circonference du cercle coupant la droite CD au point K, \& la droite CG en G. Menez sur CI les perpendiculaires KM, GL ; je dis que le sinus GL est au sinus KM, comme l'axe BF à la distance des foyers ED ;

Du point I menez \struck{la} \add{la} perpendiculaire IN sur CD \& après avoir tiré la tangente CH, \& joint CE, \struck{prolongez} \& tirez la droite EP parallele à CI sur la droite DC prolongée : \& du point C faites CO perpendiculaire à l'axe.
\note{« la » est biffé puis récrit au-dessus de la ligne.}

Ce qu'estant fait, il s'ensuit, qu'à cause que les triangles semblables CIO, \& GCL, \& les droites GC, CI, egales, que la droite CO sera egale à GL : \& à cause de l'un \& l'autre angle droit aux points M \& M, \& l'angle C commun, \& les egales CI, CK la droite IN sera aussi egale à KM.
\note{« aux points M \& M » : les deux lettres se lisent M ; la figure demande N et M. Transcrit tel qu'écrit. La page s'arrête ici ; la suite est à la vue 25.}

\page{25}
\note{Verso du feuillet de la vue 24, avec un béquet rabattu sur le bas de la page de gauche, qui masque le début d'une douzaine de lignes. Le texte n'est donné qu'une fois, à la vue 26, où le béquet est relevé. Le béquet porte, à l'encre, « 19 » en haut, et une figure d'apparence imprimée : un triangle X aux côtés hachurés, pareil à celui de la vue 21, et une construction lettrée A, B, C, D, E, F, G, I, K, L, M, N, O, P, Q, R, S, T, Y, Z, avec un arc de cercle de S à K par F et D ; le cachet de la bibliothèque y est apposé. Aucun nombre en tête de la page.}

\page{26}
\note{La même page que la vue 25, le béquet relevé vers la gauche : on n'en voit que le dos. La page poursuit la démonstration de la proposition IV.}

Donc comme CO est à IN, c'est àdire DC à DI, ou DP à DE, ainsi GL \struck{est} à KM.

\struck{Mais} Or à cause des angles ECI, ICD egaux \add{sur la} \struck{\ill{}} tangente CH, \& des paralleles CI, PE, CP sera egale à CE \& DP aux 2 droites DC, CE qui viennent des foyers DE \struck{au point} \add{à} C, qui est le point pris dans la section : c'est àdire que DP est egale à l'axe FB.
\note{« sur la » est écrit au-dessus d'un mot biffé, « à » au-dessus de « au point » biffé.}

Et partant, le sinus GL est au sinus KM comme l'axe transversant BF à la distance des foyers DE, ce qu'il falloit \struck{demonstrer} demonstrer.

Corollaire

Parceque l'axe transversant de l'Ellipse est plus grand que la distance de ses foyers \& que dans l'hyperbole il est moindre, il arrive que les loix de leurs refractions sont differentes, \& qu'elles ne peuvent servir à la mesme \ill{} dans les mesmes \& autres diaphanes, si l'on ne change leur situation.
\note{le mot entre « la mesme » et « dans » n'a pas été lu.}

\struck{La refraction de} V Proposition.

La refraction de quelconque rayon incliné sur la surface de quelconque diaphane estant donnée, descrire la figure de cette surface, qui rende tous les rayons qu'elle reçoit, \struck{ou paralleles}, \struck{\ill{}} ou qui les rassemble en un mesme point, ou qui les \uncertain{escarte} : \add{\ill{}} soit les \ill{} \uncertain{convergens} \& \uncertain{divergens}.
\note{énoncé non tenu dans une accolade. Deux mots sont soulignés, à la fin de la troisième ligne et au début de la quatrième : la lecture « convergens », « divergens » est offerte avec doute. Un mot bref, sans doute une addition, est écrit au-dessus de la ligne avant « soit ».}

Que le triangle X soit \struck{\ill{}} une piece de quelque corps diaphane comme est le cristal, \& qu'il soit transporté au triangle ABC ; le rayon soit FGI, dont l'inclination soit FID ; \& que venant de l'air FG, il se rompe en I après avoir transpercé \struck{GI} perpendiculairement GI, afin de descendre en E par IE.
\note{« en » est écrit, ici et dans tout le traité, par un signe semblable à un 3 ; il est rendu « en ».}

Soit faite DI perpendiculaire à IC, \& \struck{\ill{}} que MI face l'angle MIC egal à l'angle \uncertain{CIE}. Soit la circonference FDK \uncertain{s'estendant} par FK parallele à IC ; \& la droite KL parallele à DI, \& que le rayon IE se rencontre en L, \& soit fait tout le reste, comme dans la 2 propos.
\note{la lettre du milieu de « CIE » est couverte d'une tache.}

De plus, que IR soit perpendiculaire à DE : \& qu'ET prolongée jusques à ce qu'elle rencontre la circonference KDF au point S ; \& SV, LY soient paralleles à IC ; \& DT perpendiculaire à ES : \& \add{que le} \struck{que la section de} la droite FK couppe la droite DI soit le point Z.
\note{« que le » est écrit au-dessus de « que la section de » biffé.}

Cela fait l'on aura la raison de tous les autres rayons inclinez comme l'on voudra sur ladite surface diaphane, par le moyen de la droite DI, ou FI comparée à la droite IL, c'est àdire du rayon d'incidence, comparé au rayon rompu.

Je dis donc premierement, qu'il y a mesme raison de la droite SV à la droite FR, c'est àdire du sinus de l'angle fait par le rayon rompu \& par la perpendiculaire au sinus de l'angle fait par la mesme perpendiculaire \& par le rayon d'incidence, qu'il y a du rayon d'incidence FI au rayon rompu IL.

Car à cause des angles egaux au sommet, \& des \struck{\ill{}} angles droits en V \& Y, \& en R, comme DI, ou SI est à IL ainsi SV est à LY, ou \uncertain{HZ} ou \uncertain{FZ}.
\note{les deux dernières lettres ont une forme proche de celle du point Z défini plus haut ; plus haut, dans « la droite FR », la même lettre a la panse d'un R. La figure du béquet (vue 25) porte un point Z et n'a pas de point H.}

Je dis
\note{réclame, seule en bas à droite de la page. La suite n'est pas à la vue suivante, qui montre une bande de papier d'un autre texte (vues 27–28), mais en tête de la vue 29, qui commence par « Je dis en 2 lieu ».}

\page{27}
\note{Une bande de papier, déchirée en haut, collée par son bord gauche sur une feuille de montage grise ; la page de gauche de la vue est un papier blanc moderne. En haut à droite de la bande, « 20 », rond et appuyé, et « 17 », d'un trait plus fin. Le cachet de la bibliothèque est apposé sur la bande. La main est celle des propositions d'optique qui précèdent, mais le texte est d'une autre nature : un morceau de prière ou de méditation, sans lien visible avec la réclame « Je dis » de la vue 26. L'écriture, très rapide et chargée de ratures, est ici plus difficile que dans le traité ; beaucoup de mots ne sont donnés qu'avec doute.}

Luminaire, \struck{qui} dont les rayons \uncertain{rejallissans} par la rencontre des corps opaques, \ill{} les plus \struck{\ill{}} \ill{} \& \ill{} à leurs \uncertain{sources} \struck{\ill{}} \& nous servent de \uncertain{fondemens} \struck{pour} nous \uncertain{eslever} à \struck{\ill{}} faire \uncertain{reflechir} toutes nos pensées, \& nos desirs à l'honneur à la gloire \& à \uncertain{l'amour} de Dieu \struck{\ill{}} \uncertain{sans lequel} nous ne pouvons subsister \ill{},
\note{le premier mot, « Luminaire », ouvre la bande sans rien avant lui : la phrase commence ailleurs. Le premier mot de la quatrième ligne, bref, n'a pas été lu.}

\struck{\ill{}}
\note{une ligne entière biffée d'un trait continu.}

Je \uncertain{l'espere} donc de \ill{} faire une autre nouvelle faveur, que \ill{} \ill{}, en perfection \add{de lumiere} \struck{\ill{}} \uncertain{chose} \uncertain{aussy} qui ne \uncertain{luy} soit agreable, \& que \struck{\ill{}} \ill{} \struck{\ill{}} soit un miroir parfait qui face retour à \struck{\ill{}} \uncertain{sa} source \struck{\ill{}} d'amour toutes les inspirations qu'elle donne à nostre \struck{\ill{}} esprit, jusques à ce que cette reflexion \& cette veue par le miroir \struck{\ill{}} soit changée dans la \uncertain{claire} \ill{} du \struck{\ill{}} soleil divin, \uncertain{dans} la veue \struck{\ill{}} \& la \uncertain{jouissance} \ill{} \uncertain{bienheureuse}.
\note{« de lumiere » est écrit au-dessus de la ligne. Les mots biffés sont couverts d'un trait épais et n'ont pas été lus.}

\struck{C'est \ill{}}
\note{deux ou trois mots biffés, seuls sur la dernière ligne écrite ; le reste de la bande est blanc.}

\page{29}
\note{Recto écrit sur sa moitié supérieure seulement ; la page de gauche de la vue est le papier de montage, blanc. En haut à droite, « 21 », rond et appuyé, et « 18 », d'un trait plus fin. Le texte reprend la réclame « Je dis » laissée au bas de la vue 26 : la bande des vues 27–28 est intercalée entre les deux. Le cachet de la bibliothèque est apposé dans la marge de droite.}

Je dis en 2 lieu, qu'il y a la mesme raison de la droite ME à la droite PQ, que du rayon d'incidence FI au rompu IL. Car à cause des \struck{\ill{}} \add{des droites} paralleles FK \& ID, \& des angles FZD, IDK, \& des angles droits aux points Z \& \uncertain{R}, la droite IR sera egale à la droite FZ : \& DT \struck{\ill{}} egale à SV à cause des angles droits T \& V.
\note{« des droites » est écrit au-dessus de la ligne.}

Donc comme SV sera à FZ, ou FI à IL, ou comme le rayon d'incidence au rompu, ainsi DT sera à IR, c'est àdire ED à EI, ou EM à EN ; ou EM, qui est la distance des foyers à EN ou PQ l'axe transversant de l'hyperbole ; ou \struck{Ellipse} l'axe transversant EM de l'Ellipse, à PQ qui est la distance de ses foyers.

Je dis en 3 lieu que soit l'hyperbole, dont l'axe transversant soit à la distance de ses foyers, comme du rayon rompu IL \struck{\ill{}} au rayon d'incidence FI ; soit l'Ellipse, dont la raison de l'axe à la distance des foyers soit contraire de celle de l'hyperbole, à sçavoir la mesme que de la droite FI à la droite IL, que l'une \& l'autre de ces sections est propre pour les \struck{\ill{}} diaphanes.

C'est pourquoy l'on peut faire des Conoid. \struck{de l'hyperbole} de l'hyperbole \& des spheroides de l'Ellipse pour tous les effects qu'on desirera : car les Conoid. convexes d'un costé \& \struck{\ill{}} plats de l'autre (comme l'on fait ordinairement les verres pour les lunettes de longue veue) rendent les rayons paralleles qu'ils recevront \uncertain{convergens} ou divergens par leur costé convexe : \& rassemblent en un point \struck{les rayons} par sa surface convexe les rayons qu'ils recevront paralleles par leur costé plat \& droit.
\note{« Conoid. » : abréviation de « Conoides », deux fois. La parenthèse est de l'auteur.}

Les diaphanes spheroides feront tout le contraire, soit en recevant, soit en renvoyant les rayons. Les Conoid. concaves d'un costé \& droit de l'autre rendent les rayons paralleles par leur \struck{\ill{}} costé plat, les \uncertain{escarte} \& divise par leur costé concave ; \& si les rayons \struck{\ill{}} tendoient à s'unir, il les rendra paralleles.
\note{le texte s'arrête ici, au milieu de la page ; le bas du feuillet est blanc, à part une tache d'encre.}

\page{31}
\note{Page de droite, un feuillet plus petit monté sur onglet ; la page de gauche de la vue est une feuille de papier blanc sans écriture. En haut à droite, « 23 », rond et appuyé, et « 21 », plus fin. La main paraît la même que celle des propositions d'optique (même esperluette en grand C bouclé, même « q »), mais plus petite et plus serrée. Dans la marge de gauche, une figure à la plume : une ellipse d'axe vertical, de sommets $\zeta$ (en bas) et n (en haut), une tangente menée d'un point $\delta$ de la courbe jusqu'au point m de l'axe prolongé, deux ordonnées horizontales, celle de $\delta$ en o et celle de l en u, la lettre $\varepsilon$ entre les deux, et R, le centre. La page commence sans titre ; ce qui précède n'est pas dans ce lot. Le cachet de la bibliothèque est apposé dans le texte.}

Pour l'Ellipse \uncertain{en} \uncertain{de cette methode aux tangentes} \ill{} que \uncertain{je garde}. Soit l'Ellipse exemple $\overline{\zeta\delta\varepsilon n}$, dont l'axe soit $\overline{\zeta n}$ \& le centre R. Prenez un point quelconque $\delta$ dans sa circonference duquel il faut tirer $\overline{\delta m}$, qui touche l'Ellipse. Tirez l'appliquée $\overline{\delta o}$, \& supposez \uncertain{en premier lieu} que \uncertain{supposons} \ill{} que $\overline{o\zeta}$, qui est donnée, s'appelle B. \& $\overline{on}$ aussi donnée s'appelle G. Feignons \uncertain{que} $\overline{om}$ inconnue, s'appelle A. \uncertain{Ce qui est} \uncertain{intercepté} par $\overline{\delta m}$ la portion de l'axe \uncertain{comprise} entre le point o \& le concours de la tangente.
\note{Les noms des points sont surmontés d'un trait, ici et dans toute la page ; ils sont rendus surlignés. La lettre qui ouvre les noms, bouclée vers le bas, est lue $\zeta$ d'après la figure ; celle qui ressemble à un d bouclé, $\delta$ ; celle du milieu de la figure, $\varepsilon$. « u », « n », « m », « o » sont des lettres latines. B, G, A, E sont les lettres des espèces, écrites en capitales, G ayant la forme d'un 6 ouvert. Le premier mot suivant « Ellipse » est écrit par le signe semblable à un 3 que les feuillets d'optique emploient pour « en ». La fin de la première ligne, entre « tangentes » et « Soit », n'est lue qu'en partie.}

Puisque dans toute l'Ellipse, si on tire $\overline{lu}$ parallele à $\overline{\delta o}$, par le point u pris à discretion entre o \& n, il est clair que cette ligne $\overline{lu}$ coupe la tangente $\overline{\delta m}$, \& l'Ellipse aussi, aux deux points $\overline{\varepsilon l}$. Et parceque la \uncertain{partie} de la tangente, tous ses points \uncertain{excepté} $\delta$, \uncertain{tombent} hors de l'Ellipse, donc la ligne $\overline{lu}$ sera plus grande que la ligne $\overline{\varepsilon u}$. Il y aura donc plus grande proportion du quarré $\overline{\delta o}$ au quarré $\overline{\varepsilon u}$ que du quarré $\overline{\delta o}$ au quarré $\overline{lu}$. Or que le quarré $\overline{\delta o}$ au quarré $\overline{\varepsilon u}$, ainsi est, par la proprieté de l'Ellipse, le rectangle $\overline{\zeta on}$ au rectangle $\overline{\zeta un}$, \& que le quarré $\overline{\delta o}$ au quarré \struck{\ill{}} $\overline{lu}$, ainsi le quarré $\overline{om}$ au quarré $\overline{um}$. Donc il y a plus grande proportion du rectangle $\overline{\zeta on}$ au rectangle $\overline{\zeta un}$, que du quarré $\overline{om}$ au quarré $\overline{um}$. Feignons que $\overline{ou}$ \uncertain{prise} à discretion \uncertain{soit} egale à E.
\note{« Puisque dans toute l'Ellipse » : la lecture du troisième mot est incertaine. Le mot biffé avant « $\overline{lu}$ » est court.}

Le rectangle $\overline{\zeta on}$ se \uncertain{nommera} B in G. Le rectangle $\overline{\zeta un}$ se \uncertain{nommera} B in G $-$ B in E $+$ G in E $-$ Eq. Le quarré $\overline{om}$ se \uncertain{nommera} Aq. Et le quarré $\overline{um}$ se \uncertain{nommera} Aq $+$ Eq $-$ A in E bis. Il y a donc plus grande \struck{\ill{}} proportion de B in G à B in G $-$ B in E $+$ G in E $-$ Eq, que de Aq à Aq $+$ Eq $-$ A in E bis.
\note{« se nommera » : le verbe, écrit vite et chaque fois de la même façon, est offert avec doute ; « Aq » et « Eq » sont les quarrés de A et de E, « bis » le double, « in » le produit, dans la notation des espèces.}

Par consequent si l'on multiplie la premiere par la derniere, \& la 2 par la 3. B in G in Aq $+$ B in G in Eq $-$ B in G in A in E bis, qui est le produit du premier par le dernier, sera plus grand que B in G in Aq \struck{$+$ B in G in Eq $-$ B in G in A in E bis qui est le produit \ill{}} $-$ B in E in Aq $+$ G in E in Aq $-$ Aq in Eq. \struck{derniere, \& la plus grande}
\note{le passage biffé répète par erreur le premier produit ; il est barré d'un trait continu qui court sur la fin de la ligne et le début de la suivante.}

Il faut donc suivant la \uncertain{susdite} methode, comparer ces 2 produits par adæqualité. Ostant ce qui est commun de part \& d'autre, \& divisant le reste par E, il restera d'un costé B in G in E $-$ B in G in A bis, \& de l'autre $-$ B in Aq $+$ G in Aq $-$ Aq in E. Effaçant les termes \uncertain{qui sont affectez} de la ligne E, il restera d'un costé $-$ B in G in A bis, \& de l'autre $-$ B in Aq $+$ G in Aq, lesquels il faut maintenant, suivant la \uncertain{susdite} methode, \uncertain{egaler tout alors}, \& divisant le tout par A, \& transposant \uncertain{d'un costé les termes contraires}, B in A $-$ G in A egal à B in G bis.
\note{« adæqualité » : le mot, écrit en un seul trait, est lu avec la ligature æ. Le cachet de la bibliothèque recouvre le premier « B » de « B in G in E », et une tache d'encre la lettre qui suit « A » dans le même terme.}

\uncertain{Voila donc} que cette solution de la question \uncertain{est la mesme que} celle d'Apollonius, car \uncertain{selon ma construction}, pour trouver la tangente, il faut faire que B $-$ G est à G, ainsi B bis à A, c'est àdire que $\overline{\zeta o}$ moins $\overline{on}$ à $\overline{on}$, ainsi $\overline{\zeta o}$ bis à $\overline{om}$. Et par celle d'Apollonius, il faut faire que $\overline{\zeta o}$ à $\overline{on}$, ainsi $\overline{\zeta m}$ à $\overline{mn}$. \uncertain{Or} ces 2 constructions, qui il paroist, \uncertain{reviennent} à la mesme.
\note{« $\overline{\zeta m}$ à $\overline{mn}$ » : les deux lettres, liées, peuvent se lire aussi u et n.}

Je pourrois adiouster plusieurs autres exemples tant du 1\textsuperscript{er} que du 2 cas de ma methode, mais ceux cy suffisent pour faire assez voir qu'elle est generale, \& qu'elle ne trompe jamais.

Je n'adiouste point non plus la demonstration de la regle, ni plusieurs autres usages \uncertain{qui en font encore} paroistre la perfection, \uncertain{comme} de l'invention des centres de gravité, dont j'ay envoyé \uncertain{les exemples} à M\textsuperscript{r} Roberval, \& plusieurs autres, \uncertain{me contentant} de proposer la suivante question à soudre à ceux qui voudront l'entreprendre.
\note{la page s'arrête là ; la question annoncée n'est pas écrite sur ce feuillet. Le verso (vue 32, page de gauche) est blanc.}

\page{32}
\note{Un bifeuillet plié, posé ouvert sur la page de gauche de la vue : sa page de gauche est écrite, sa page de droite blanche porte en haut à droite « 24 », rond et appuyé. La main est celle de la vue 31, petite et rapide, mais l'écriture est ici plus lâche et plus difficile ; plusieurs lignes ne sont lues qu'en partie. Le texte commence en haut de la page, sans titre. Le cachet de la bibliothèque n'y est pas apposé.}

M\textsuperscript{r} \uncertain{Fermat} dans sa derniere let. de Juillet 1638 dit que M\textsuperscript{r} de \ill{} s'est trompé en sa methode de \ill{}, puisqu'il a creu qu'il falloit bailler 2 noms à la ligne que j'appelle B, dont l'un est $\dfrac{\text{B in A} - \text{in E}}{\text{A}}$, l'autre $\dfrac{\text{B in A} + \text{B in E}}{\text{A}}$, \& que l'autre luy doit estre baillé \uncertain{selon} la proprieté specifique de la courbe, ce qui reussit \uncertain{en la parabole} \& autres questions \uncertain{aysées}, \uncertain{mais aux} difficiles, il est impossible d'y venir à bout. \ill{} \ill{} qu'on peut \uncertain{bailler} \uncertain{un nom} à la ligne \uncertain{selon} la proprieté specifique de la courbe, \uncertain{ce qui se peut trouver} \uncertain{par ma methode}, \ill{} \ill{} \ill{}, \ill{} \ill{}.
\note{« let. » : abréviation de « lettre ». Le nom qui suit « M\textsuperscript{r} de », au début de la ligne, commence par un C et n'a pas été lu ; plus bas, la même page nomme « M\textsuperscript{r} de Cartes » (voir ci-dessous). Dans la première fraction, « in E » est écrit sans lettre devant, tel quel. Les deux fractions sont écrites comme ici, avec un trait et le A au-dessous. Les cinq dernières lignes de ce paragraphe sont écrites très vite et ne se lisent que par endroits ; la lecture proposée est fragmentaire.}

Je dis de plus que \uncertain{par ma methode} \ill{} \& plus \uncertain{aysement} que \uncertain{la sienne}, car \uncertain{soit} luy proposée la courbe de laquelle la proprieté specifique soit exprimée par cette egalité
\[
x^{10} + g\,x^{9} + g^{2}x^{8} + g^{3}x^{7} + g^{4}x^{6} + g^{5}x^{5} + g^{6}x^{4} + g^{7}x^{3} + g^{8}x^{2} + g^{9}x
\]
\[
\propto
\]
\[
y^{10} - f\,y^{9} - f^{2}y^{8} - f^{3}y^{7} - f^{4}y^{6} - f^{5}y^{5} - f^{6}y^{4}\ f^{7}y^{3} - f^{8}y^{2} - f^{9}y
\]
\note{les deux membres sont écrits sur deux lignes, le premier se poursuivant sur une troisième par « $+ g^{9}x$ ». Entre les deux, sous le premier, un signe isolé en forme de « \& » renversé, rendu $\propto$ : c'est vraisemblablement le signe d'égalité de la Géométrie, mais la lecture est une inférence. La lettre rendue $g$ a la forme d'un 6 ouvert, la même qu'à la vue 31 où elle a été lue G ; la lecture b n'est pas exclue. Devant $f^{7}y^{3}$, aucun signe n'est écrit. L'exposant de $g^{8}$ est récrit.}

Je dis que \uncertain{par ma methode} on fera la solution sans nulle difficulté, car il ne faudra que \uncertain{traiter} la question \uncertain{à l'ordinaire} par adæqualité, en prenant $\dfrac{\text{B in A} - \text{B in E}}{\text{A}}$ pour x au lieu de x, \& D $-$ E pour y au lieu de y, \& jamais une seule \uncertain{asymmetrie} ne s'y \uncertain{rencontrera}, ce qui fait voir la facilité \& perfection de cette methode, \& \uncertain{qu'elle n'a point les} difficultez qu'a celle de M\textsuperscript{r} de Cartes, de s'embarrasser dans les \uncertain{asymmetries}, car lorsqu'il se sert de sa methode, où il s'efforce de rendre \uncertain{egales} les \uncertain{racines} par le moyen de la multiplication (page 344. l. 3.) il \uncertain{se void fort empesché}.
\note{« page 344. l. 3. » est net et entre parenthèses : un renvoi à une page et à une ligne d'un imprimé, qui n'est pas nommé.}

\uncertain{Or} puisque sa solution a \uncertain{resolu} tout ce que je luy ay proposé \uncertain{touchant} \ill{}
\note{la ligne se termine par un long trait de plume ; la suite, s'il y en a une, n'est pas sur la page.}

\note{Au-dessous, une figure à la plume : six demi-cercles de plus en plus grands, tous issus d'un même point A à gauche et posés sur une même droite, qui porte à la suite les lettres B, C, D, E, F, G ; une perpendiculaire élevée d'un point y de cette droite les coupe aux points marqués, de bas en haut, o, u, e, s, n, m. La lettre A et la lettre y sont écrites sous le pied de la figure et sont en partie effacées.}

\begin{enumerate}
\item Sunt 6 semicirculi AoB, AuC, A\uncertain{e}D, a\uncertain{s}E, anF, amG, super eadem recta AB. Ita secare in y, ut applicatâ perpendiculari yom occurrente semicirculis datis, summa rectarum yo, yu, ye, ys, yn, ym, sit aequalis rectae datae.
\item Secetur AB in y ut summa rectarum yo, yu, ye, ys, yn, ym sit \struck{aequalis rectae} \struck{datae} \uncertain{minima vel maxima}.
\item Dati coni scaleni superficiei invenire aequale planum.
\item Invenire triangula rectangula in numeris, quorum area aequetur quadrato, aut demonstrare impossibilitatem quaestionis.
\end{enumerate}
\note{les quatre questions, en latin, sont numérotées 1 à 4 dans la marge de gauche. Les lettres des demi-cercles sont données d'après la figure, les points e et s y étant peu nets. Dans la deuxième, « aequalis rectae » est biffé à la fin d'une ligne et « datae » au début de la suivante ; ce qui suit est écrit vite. Rien ne dit sur la page si ces questions sont celles que la fin de la vue 31 annonce (« la suivante question ») : le rapprochement n'est qu'une possibilité.}

\page{33}
\note{Page de droite, un feuillet monté sur onglet ; la page de gauche de la vue est une feuille de papier blanc sans écriture, pliée. En haut à droite, « 25 », rond et appuyé, souligné d'un trait, et « 12 » d'un trait plus fin. La main paraît celle des pages d'optique et de la vue 31, écrite vite, avec de nombreuses ratures et additions interlinéaires. Le texte commence en haut de page par un signe de renvoi (une croix surmontée de trois points) et « Car cela est », sans titre : il reprend une discussion commencée ailleurs. Il porte sur la vitesse de la lumière, comparée à celle des boulets, des flèches et des étoiles. Dans la marge de gauche, écrite en long de bas en haut sur trois lignes, une longue addition de la même main. Le cachet de la bibliothèque est apposé au milieu de la page.}

\uncertain{Car cela est}

Puisque \struck{la lumiere} \uncertain{cette} la \uncertain{meilleure} de la \uncertain{perfection} \& il \struck{\ill{}} \add{nous rend dans ses \uncertain{rapports} \ill{}} \ill{} sa creature, ou du moins \struck{\ill{}} \add{\ill{}} \uncertain{respec.} \&c. Quant à sa creature il est facile que nous \uncertain{prenions} la \uncertain{premiere} \uncertain{sans} une idée plus intelligible que \uncertain{sera} celle du mouvement le plus \uncertain{prompt} que nous puissions \struck{\ill{}} imaginer.
\note{les deux premières lignes sont chargées de ratures et d'additions au-dessus de la ligne, en partie illisibles ; la lecture proposée n'en donne que l'ossature.}

Car si nous considerons la communication qu'a le soleil de ses rayons depuis l'orient jusques à l'occident, lors que nous le voyons dans le soleil, \& dans les estoilles dont elle vient jusques à nous \& qu'elle retourne dans sa source, \& dans tous les lieux du monde, suivant les corps opaques qui la font reflechir, \& \uncertain{pour} faire \uncertain{ensuite} autre \uncertain{chose} \& toute sa course, qu'elle fait presque en \uncertain{instant} sans nous \uncertain{assurer} qu'elle n'a point de vitesse \uncertain{determinée} qui l'egale.

Ce qu'il est aysé de \uncertain{prouver} par l'exemple de \struck{\ill{}} plus grande vitesse, dont je fais \uncertain{juger} \uncertain{tout} \uncertain{naturel}, affin que ceux qui ne sont pas accoustumez à considerer une si grande vitesse que celle de la lumiere, \& jugent de sa facilité.

Je \uncertain{commence} donc par la vitesse d'un cheval, \struck{\ill{}} \uncertain{lequel} le plus viste qui se puisse \uncertain{trouver} \uncertain{qui} \uncertain{courant} la course fait seulement 21 toises en trois secondes, c'est àdire 7 toises ou 42 pieds dans le temps d'une seconde minute, de sorte que s'il pouvoit courir de la mesme vitesse la dixiesme partie d'une heure ou 6 minutes il feroit une lieue, \& par consequent dix lieues dans \uncertain{une} heure.

Et parceque l'on tient \struck{\ill{}} \add{quelques} \uncertain{levriers} ou \struck{\ill{}} \add{\ill{}} \uncertain{chiens} vont deux fois plus viste que le cheval \struck{\ill{}} \struck{\ill{}}, \& qu'on tient de \uncertain{plus} que l'oyseau qui vole le plus viste va deux fois aussi viste que le cheval, il s'ensuit que l'oyseau fait 28 toises dans une seconde minute, ou 40 lieues dans une heure, ce qui \add{me} \uncertain{semble} probable, puisque la fleche tirée horizontalement ne fait que 30 toises \struck{dans deux secondes}, \& que \uncertain{communement} elle va plus viste que l'oyseau, qui sans doute seroit frappé de la fleche tirée avec un arc de bois, ou du moins du matras tiré avec une arbaleste d'acier, qui ne fait que 30 toises dans une seconde.
\note{les chiffres 21, 7, 42, 28, 40 et 30 sont nets ; le raisonnement qui les relie est donné tel qu'il se lit, sans être corrigé. Le passage sur les chiens et l'oiseau est surchargé d'additions interlinéaires ; « levriers » et « chiens » sont offerts avec doute. « matras » : le trait d'arbalète.}

Or l'experience \uncertain{m'enseigne} que la bale d'arquebuse va \struck{\ill{}} 4 fois plus viste que \uncertain{l'arc} \uncertain{susdit}, car elle fait 120 toises dans une seconde : \& parcequ'il semble que nous n'ayons point \add{de vitesse} plus grande que celle qui sort du \uncertain{mousquet} \uncertain{ou} canon, \struck{\ill{}} \add{\ill{}} \add{laquelle} \uncertain{ne cede} \add{\uncertain{peut estre}} pas à la vitesse \struck{du foudre} de la poudre ou du feu qui accompagne le tonnerre, nous la pouvons comparer à la vitesse du mouvement de la terre, si c'est elle qui fait le tour \& le circuit de 24 heures, car puisqu'elle a 7200 lieues en \uncertain{circonference}, elle fait 300 lieues à chaque heure ou 5 lieues à chaque minute qui reviennent à 208 toises \struck{\ill{}} \uncertain{à peu pres}, pour chaque seconde.

Cette vitesse est bien grande ; mais celle des estoilles \struck{\ill{}} qui font le tour du monde en 24 heures, car \uncertain{nous tenons} \struck{\ill{}} est beaucoup plus grande si elles \uncertain{sont} au firmament 14000 fois plus grand que celle de la circonference de la terre. D'où il arrive qu'il faut multiplier la mesme \uncertain{nombre} de sa vitesse, \& partant que les estoilles qui font leur \uncertain{cours} dans une \uncertain{si} grande estendue de leur ciel en 24 heures font plus de \uncertain{mille} \uncertain{cinquante} lieues dans une seconde \uncertain{quand} nous \uncertain{prenons} la moindre grandeur que les \uncertain{Astronomes} luy donnent au 8 ciel.
\note{« 7200 », « 300 », « 5 », « 208 », « 14000 », « 24 » sont nets. Le chiffre des lieues parcourues par les étoiles en une seconde est écrit en toutes lettres et n'est pas sûr. La page s'arrête ici.}

\marginal{\uncertain{Mais} \uncertain{Roberval} \ill{} \struck{\ill{}} \add{\ill{}} \ill{} dans la \uncertain{proportion} que \uncertain{peut faire} \uncertain{communiquer} \ill{} à la \ill{} \uncertain{celuy} qui a tant d'\uncertain{honneur} qu'il suffit de l'\uncertain{avoir} pour \ill{}, c'est la \uncertain{perfection} par laquelle \ill{} \ill{}. Or puisque \uncertain{ce} \ill{} est d'un \uncertain{estre} \uncertain{plus} de sa \ill{} \uncertain{en} \uncertain{nous} \ill{} \ill{} la contemplation de cette \ill{} \ill{} \uncertain{divine} \ill{} \uncertain{pour} plusieurs \ill{} de \uncertain{perfection} \ill{} \ill{} \& \ill{} \ill{} a \ill{} plus \ill{}, \ill{} \ill{}.}
\note{addition écrite en long dans la marge de gauche, de bas en haut, sur trois lignes, de la même main et de la même encre, sans signe de renvoi visible ; on la donne ici, après le texte. Elle est serrée contre l'onglet et presque entièrement illisible à cette résolution ; seuls quelques mots se lisent. Le premier mot abrégé, « M\textsuperscript{r} Rob. » ou « Mais », n'est pas sûr.}

\page{34}
\note{Ouverture : à gauche, le verso du feuillet de la vue 33 ; à droite, le recto du feuillet suivant. Les deux pages sont de la même main rapide, chargées de ratures et d'additions entre les lignes, et beaucoup de mots ne se lisent qu'avec doute. La page de gauche poursuit la méditation sur la vitesse de la lumière de la vue 33 et la tourne vers les corps glorieux ; la page de droite (ci-dessous) continue sur les « douaires » des corps glorieux, puis sur la lumière qui reproduit partout l'image de son principe, et passe enfin à la portée de la vue sur le globe terrestre. Dans la marge de gauche de la page de gauche, une petite figure : un cercle lettré A, B, C, D, et ses deux diamètres. En haut à droite de la page de droite, « 26 », rond et appuyé, et « 13 », plus fin.}

Car si nous prenons celle qui luy est necessaire lors que \uncertain{l'on suppose} les 7 planettes ordinaires ne doit estre pris que pour un point à son egard, nulle raison \uncertain{n'empesche} que nous ne le supposions cent milliers de fois plus grand ; de sorte que son diametre \uncertain{seroit} 100000000 fois celuy de la terre, ou \uncertain{pourroit}, \add{s'il ne falloit}, \uncertain{le nombre} des milliers de fois plus grand que le precedent pour \ill{} sa grandeur.
\note{la construction de cette phrase n'est pas assurée ; « 100000000 » est net.}

Ce qui augmentera sa vitesse en mesme proportion, \& luy fera faire plus de chemin dans une seconde, qu'elle n'en a fait depuis qu'elle a esté creée jusques à present.

Or bien que cette vitesse semble incomparable, elle ne l'est pas si on la compare à celle d'une \ill{}, quand on l'applique à la vitesse de la lumiere laquelle semble estre la plus grande de toutes les possibles \& surpasser toutes les autres, voire \uncertain{la surpasse infiniment} \struck{\ill{}} \uncertain{en} toutes sortes de façons.

\struck{Cette vitesse de la lumiere se \ill{} plus \ill{}} par la figure du cercle ABCD, \struck{\ill{}} supposé qu'il soit egal au plus grand cercle de l'univers \& qu'il comprenne tout l'espace creé, de sorte que le diametre DB soit \struck{infiniment} 100000000 fois \struck{autant} plus grand que l'espace \struck{\ill{}} du pole arctique à l'antarctique, il est clair que si le soleil, ou une estoille, ou tel autre luminaire qu'on voudra est imaginé en tel point, comme au point D, d'où il envoye sa lumiere jusques à B, \& que s'il trouve \struck{\ill{}} un corps reflechissant de B en A, \& d'A en C, \& de C en D, \& ainsi à l'infini, \uncertain{nous verrons} tout ce chemin DBACD, \& un million d'autres semblables, en un temps si court \& si imperceptible que nous ne le pouvons mieux signifier que par un instant ou un moment.
\note{« DBACD » est net.}

Et si \add{nous} \uncertain{l'on} considerons le rayonnement d'un point lumineux en D, ou en B, ou en tel autre point de la circonference de ce cercle, ou de la surface de ce globe \struck{\ill{}} \add{\uncertain{tournant} \uncertain{tout} \uncertain{autour}}, qu'il remplit dans un moment toute la solidité de cet espace, sans qu'il se puisse trouver aucun point où il ne pousse un rayon ; de sorte que cette sphere d'espace est toute remplie de lumiere.

Ne pouvons nous pas conclure que \uncertain{le} corps glorieux a le don, ou la \uncertain{vertu} d'agilité, par laquelle il peut aller aussi viste que la lumiere ? \& \struck{\ill{}} la \uncertain{clarté} ou la lumiere ne luy peut elle pas aussi bien estre attribuée ? \add{qui le \uncertain{rend}} \struck{\ill{}} qu'il peut \struck{\ill{}} tellement \uncertain{\ill{}} qu'il fera paroistre ou disparoistre sa lumiere quand il luy plaira.
\note{deux questions, chacune fermée par un point d'interrogation ; le verbe avant « qu'il fera paroistre » est surchargé.}

Car 2 \struck{\ill{}} proprietez des corps \uncertain{bienheureux} estant posées, il est evident \struck{\ill{}} qu'il peut \struck{\ill{}} tellement remplir le ciel \& la terre, à nostre egard, que nous le verrons \& le toucherons par tout. \add{\ill{}} nous pouvons \uncertain{imaginer} qu'il \add{\ill{}} \struck{\ill{}} \add{\uncertain{sa vitesse}} \struck{\ill{} \ill{} sa puissance \ill{}} ; car pour oster le moment qui semble incomparable, \uncertain{puisqu'il} faut \uncertain{un moment} ou \ill{} pour \uncertain{comprendre} un si grand \uncertain{univers}, \uncertain{supposons} \struck{\ill{}} le temps d'un \add{\uncertain{battement}} \ill{} qui est la soixantiesme partie d'une seconde minute, \& que \struck{\ill{}} pendant ce temps le corps glorieux \struck{\ill{}} \add{passe} par toutes les lignes \uncertain{de l'espace} total de la sphere ABCD, par lesquelles le rayon du luminaire D passe dans \uncertain{un} moment, \uncertain{il} peut \uncertain{remplir} le tout dudit espace, \uncertain{estre} \uncertain{touché} par le corps glorieux, puisqu'il n'y a nul espace où il ne soit dans le temps d'une tierce.
\note{« tierce » : la soixantième partie d'une seconde, comme la phrase le dit ; le dernier mot est écrit en abrégé et la lecture en est offerte par le sens. Toute la fin de cette page est surchargée ; la lecture n'en donne que l'ossature.}


\note{Page de droite de la vue 34.}

Et si nous voulons supposer qu'il ait le douaire ou la proprieté d'impassibilité, il \struck{\ill{}} \uncertain{les} \uncertain{pourra} \ill{} tous les autres corps \struck{\ill{}} glorieux qu'il rencontrera s'il \uncertain{jouit} de l'autre douaire de subtilité, par laquelle il passera à travers tout autre \struck{\ill{}} corps \uncertain{comme} le rayon \uncertain{passe} \uncertain{par} le verre ou le cristal, sans qu'il \struck{\ill{}} \uncertain{l'empesche}.
\note{« douaire » : les dons du corps glorieux (impassibilité, subtilité, agilité, clarté), dans la langue des théologiens ; le mot revient plus bas.}

Enfin si nous supposons que ces proprietez soient tellement en la disposition des \uncertain{bienheureux} qu'ils les puissent oster ou les \uncertain{reprendre} selon \struck{\ill{}} \uncertain{leur volonté}, qu'ils les \uncertain{puissent} \uncertain{donner} comme ils voudront, \uncertain{il s'ensuit} qu'un corps glorieux peut \uncertain{estre} lumineux à toutes les \uncertain{heures} ; car lors qu'il \struck{\ill{}} voudra, allant de la grande vitesse dont nous avons parlé, il se voudra seulement rendre lumineux dans \uncertain{la} partie du monde où il luy plaira, \& il sera \uncertain{luy mesme} \add{\uncertain{donnera}} \struck{\ill{}} plus ou moins de clarté, s'il a la puissance de \uncertain{la modifier} \struck{\ill{}} comme il luy plaira.

D'où il est aysé de conclure beaucoup d'autres choses en faveur de la \uncertain{perfection} du corps glorifié \uncertain{nostre} \uncertain{felicité}.

Mais avant que de quitter la lumiere, nous pouvons remarquer qu'elle sert particulierement à produire l'image de son principe dans tous les lieux du monde, comme si elle estoit à la \uncertain{representer} \& reproduire par tout, \& à la representer d'une immensité pour remplir \struck{tout} le monde \uncertain{entier}.

Car si elle trouvoit tous les corps qu'elle frappe, polis comme \struck{\ill{}} doit estre dans la \uncertain{surface} d'un miroir parfait, nous ne verrions rien que l'image de son principe par exemple, que le soleil, \struck{\ill{}} si la lumiere en \uncertain{parloit}, de sorte que chaque miroir qu'il frappe \struck{\ill{}} \uncertain{luy} \uncertain{rendroit} \ill{} \struck{\ill{}} l'impression : de sorte que les couleurs ne sont autre chose que la lumiere diversement reflechie par une infinité de petits miroirs de chaque corps \add{qui ne sont pas assez polis}, qui \uncertain{feroient} \uncertain{voir} \uncertain{continuellement} le corps solide du soleil, ou des autres principes de la lumiere.
\note{l'image du miroir parfait qui renvoie la lumière à sa source est déjà celle de la bande de la vue 27.}

Et s'il nous est permis de nous \uncertain{eslever} plus haut, nous pouvons dire que toutes les creatures \uncertain{n'ont} point d'autre but principal que de nous representer leur principe, \& que toutes leurs actions \uncertain{tendent} à \uncertain{publier} \uncertain{sa} \uncertain{bonté}, sa puissance \& sa bonté \uncertain{infinie} par tout l'espace \uncertain{creé}, qu'elle \uncertain{remplisse}, soit fini ou infini.

Or je \uncertain{veux} achever cette lettre par la consideration d'un seul rayon, qui par l'effet nous pouvons \uncertain{recognoistre}, quand il nous plaira, \uncertain{trouver} la grandeur de toute la terre, \& par lequel un homme posté sur \add{\ill{}} \uncertain{sa} \uncertain{tour} \uncertain{de} la lune ou d'une \struck{autre} estoille, peut reconnoistre, par la \uncertain{veue} seulement, qu'il est plus sur la terre \ill{} \uncertain{que} \uncertain{luy} \ill{} qu'il \uncertain{visite} des \uncertain{planettes}, \& d'autres corps semblables aux nostres dans d'autres globes.
\note{« achever cette lettre » : lecture offerte avec doute, mais elle ferait de ces feuillets une lettre. Tout le paragraphe est écrit très vite ; la construction n'est pas assurée.}

Soit donc le globe terrestre ABC, \& \struck{\ill{}} \add{que nostre oeil soit eslevé} de 6 pieds \struck{\ill{}} au dessus de sa surface, par exemple en D, lors que nous \uncertain{regardons} une \uncertain{allée} \struck{d'une lieue de long}, ou un estang ou riviere d'une lieue de long, ou moins. \struck{\ill{}} Si nous \uncertain{voulons} \uncertain{sçavoir} le lieu où nous commençons à ne voir plus rien sur la terre ou sur l'eau, soit la \uncertain{lumiere} \uncertain{partant} d'une chandelle ou de quelque autre corps \uncertain{lumineux}, qui est de 6 pieds \uncertain{eslevé} au lieu où nous commençons de voir la surface de ce globe, par exemple au point F, cet espace, \struck{\ill{}} c'est àdire la ligne AF, \struck{qui n'est pas sensiblement} \uncertain{peut} \uncertain{estre} supposée egale au rayon DF, parceque leur difference ne nous est pas sensible, veu d'autre \uncertain{part} le diametre AB ;
\note{dans la marge de gauche, en regard, une figure : un cercle de diamètre vertical AB, un point E au-dessus de A sur le diamètre prolongé, une droite menée de E et tangente au cercle en C, une autre courte droite de D à F, F sur le cercle près de A. Le texte nomme D l'œil, F et C des points de la surface ; la figure porte les lettres E, D, A, F, C, B.}

Ce que je prouve parceque le quarré du rayon DF ou de la distance AF, que nous prenons ici pour la mesme chose, est egal au rectangle de BDA, de sorte qu'ayant quarré \struck{\ill{}} la distance AF \ill{} \ill{}
\note{la page s'arrête ici en pleine phrase ; la suite est au verso de ce feuillet, page de gauche de la vue 35 (« Car puisque AD la hauteur de nostre oeil… »).}

\page{35}
\note{Ouverture : à gauche, le verso du feuillet paginé 26 (vue 34, page de droite) ; à droite, le recto du feuillet suivant, avec un béquet rabattu sur son milieu. La page de gauche achève la lettre des vues 33–35 ; elle est donnée ici. La page de droite est donnée à la vue 36, où le béquet est relevé. Dans la marge de gauche de la page de gauche, en haut, la figure de la vue 34 reprise en petit : un cercle de diamètre vertical, lettres D, A, G, C, B. Le béquet de la page de droite porte une figure à la plume : un trapèze allongé, la base supérieure lettrée A, K, E, B, deux diagonales, un cercle inscrit, et au-dessous des droites convergeant vers un point I, lettres C, F, D ; le cachet de la bibliothèque y est apposé.}

Car puisque \struck{AD} la hauteur de nostre oeil est supposée d'une toise, \& la distance \add{\ill{}} \struck{\ill{}} AF ou DF d'une lieue ; \& que la lieue a 2500 toises, \uncertain{veu} qu'elle \uncertain{a} la lieue depuis la \uncertain{Bastille} jusques au Chasteau \uncertain{du} Bois de Vincennes, il s'ensuit que le rectangle BAD contiendra 6250000 toises, à sçavoir autant qu'une lieue en quarré : \& parceque ce rectangle n'aura pas plus de largeur que d'une toise, \& qu'elle ne multiplie point, le diametre BA aura six millions deux cens cinquante mille toises, \uncertain{moins une}, qui ne doit pas nous faire changer de nombre, \struck{\ill{}} \add{dont} si nous \struck{\ill{}} \add{prenons} \struck{\ill{}} la moitié AG, le rayon de la terre sera de 3125000 toises, lesquelles \uncertain{font} 1250 lieues.
\note{la lieue est prise de 2500 toises ; le calcul est juste : $2500^2 = 6\,250\,000$, et la moitié, $3\,125\,000$ toises, fait 1250 lieues. Le lieu de la mesure, « depuis la Bastille jusques au Chasteau du Bois de Vincennes », n'est lu qu'avec doute.}

D'où il est aysé de trouver la grandeur du tour de la terre \struck{\ill{}} \add{\ill{}} \uncertain{plus} \uncertain{grande}, puisqu'il \struck{\ill{}} est à BA, comme 22 à 7 : de sorte qu'il ne faut que prendre 3 fois AB avec une septiesme partie d'iceluy AB, \struck{qui fait \ill{}} ou faire que comme 7 est à 22, ainsi soit le diametre AB de 2500 à 7857 $\frac{1}{7}$ lieues, qui donnent \struck{\ill{}} le tour de la terre.
\note{$2500 \times 22/7 = 7857\frac{1}{7}$ : le calcul est juste.}

Et parceque le soleil est plus grand, le rayon visuel DF sera plus grand, comme il sera \uncertain{moindre} \ill{} \uncertain{pour ceux} qui ne seront pas sur la terre. \struck{\ill{}}

Je sçais que l'observation sera d'autant plus exacte que l'oeil sera plus eslevé sur la terre, à raison que l'angle de la vision sera d'autant plus aigu ; mais il suffit que je vous aye fait voir la subtilité de la lumiere, à qui le rayon DF est essentiel, \& que vous puissiez prendre les lieux avantageux \struck{\ill{}} comme sont les plus hautes tours, \& les montagnes, quoy que la grande distance qu'il y a de leurs sommets soit bien difficile à mesurer, pour avoir la longueur du rayon optique : \& que les \uncertain{moindres} \struck{\ill{}} \uncertain{servent} aussi pour \uncertain{sçavoir} les distances, \struck{\ill{}} \& que la moindre hauteur ne \uncertain{soit} point plus \uncertain{basse} que 3 pieds, ou \ill{} \uncertain{environ} de la hauteur, \& qu'on peut baisser l'oeil D jusques à ce qu'il \uncertain{rase} la riviere, ou les estangs, \uncertain{pendant} par un lieu de longueur.

Je laisse mille considerations \struck{\ill{}} utiles \& agreables de peur de vous ennuyer par une si longue lettre, qui suffit pour vous persuader que la science de l'optique merite d'estre prisée, \& qu'elle fournit autant ou plus de sujets pour nous eslever à Dieu, qu'aucune autre science, ce qui \uncertain{me} persuade que vous la cherirez, \& que vous ne serez pas marry que j'aye employé quelque temps à la rendre françoise pour ceux qui n'entendent point d'autre langue.
\note{la lettre finit ici, sans formule ni signature ; le bas de la page est blanc. Elle se donne pour la présentation d'une optique mise en français (« à la rendre françoise ») ; ni l'auteur, ni le destinataire, ni l'ouvrage ne sont nommés sur ces pages. Le rapprochement avec les propositions d'optique des vues 21–29 est possible, mais la lettre ne le fait pas.}

\page{36}
\note{La même ouverture que la vue 35, le béquet de la page de droite relevé vers la gauche : on n'en voit que le dos, avec un petit papillon blanc collé dessus. La page de gauche est celle de la vue 35, donnée là. La page de droite, recto d'un nouveau feuillet, commence sous le titre « Propositions » ; elle est de la même main, plus posée. En haut à droite, « 40 », d'un trait fin, et plus bas « 28 », rond et appuyé. C'est la première fois dans ce lot que le petit nombre est au-dessus du grand, et que le petit dépasse le grand : voir l'en-tête. Le cachet de la bibliothèque est apposé au milieu de la page.}

Propositions
\note{titre, en haut à droite de la page.}

Les rayons tombans d'un corps lucide sur un \struck{\ill{}} plan diaphane different en transparence du milieu par où ils passent pour y tomber, se rompent tous de mesme façon, \struck{\ill{}} si leur cheute fait des angles egaux, \& s'ils font des angles differens ils se rompent \struck{\ill{}} differemment.
\note{énoncé tenu dans une accolade.}

Soit le corps lucide AB, qui envoye ses rayons \struck{\ill{}} passans par l'air ABCD sur le plan CD du cristal CFDGC, dont la premiere surface CD est droite, \& plate, \& la seconde CGD est courbe, soit que sa courbeure soit spherique ou hyperbolique, car cela n'est pas maintenant en question.

L'experience enseigne que tous les rayons AKE, \& ceux qui tombent à plomb sur la glace CD du diaphane CFD, passent à travers sans se rompre, \& \uncertain{les} autres se rompent sur cette surface CD comme l'on void aux rayons BM, \& AO, car au mesme moment que \struck{\ill{}} le rayon BM \struck{\ill{}} rencontre au point M de la surface FD, \uncertain{à cause de} l'inclinaison du rayon BM, il se \uncertain{plie} pour aller au point N, où il se \struck{\ill{}} rompt encore à cause de la ligne courbe GD, qu'il rencontre au point N.

Les rayons BC, \& AD se rompront encore davantage, parceque les angles qu'ils font avec la ligne CD sont plus \add{differens} \struck{\ill{}} des angles droits que font les rayons qui tombent à plomb.

Mais il est plus difficile de trouver combien chaque rayon se rompt plus ou moins l'un que l'autre ; nous \uncertain{sçavons} seulement icy que les rayons de chaque costé qui tombent à pareils angles se rompent en telle sorte qu'ils se doivent rencontrer dans un mesme point ; par exemple \struck{\ill{}} si le rayon AC se rompt par la rencontre de la convexité du verre CL, en I, le rayon BD se rompra de mesme de D en I, où il se rencontrera.

Et si le rayon BM \add{du costé droit} se rompt tellement aux points M \& N qu'il aille en I, le rayon du costé gauche \uncertain{sera} \uncertain{comme} DA en O. Et si le rayon B se rompoit de CDH, son \uncertain{correlatif} AD l'iroit aussi trouver en H, \& ainsi de tous les autres ; de sorte que si apres que tous les rayons qui tombent à plomb sur CD, \struck{\ill{}} \add{par exemple} les 7 qui sont icy paralleles au droit BD, ont \struck{\ill{}} traversé \struck{\& qu'ils passent \ill{}} CD sans se rompre, \struck{\ill{}} rencontroient une surface courbe, comme CGD, sur laquelle ils tombassent tous à angles egaux, ils se ramasseroient tous dans un mesme point, comme il arrive en cette figure qu'ils se \uncertain{reunissent} \add{tous} au point I.
\note{« le rayon B se rompoit de CDH » est donné tel qu'il se lit ; la lettre qui suit B n'est pas écrite.}

Nous \struck{\ill{}} monstrerons apres que cette ligne \add{circulaire} \uncertain{doit} estre hyperbolique \& que la \struck{\ill{}} spherique \uncertain{s'approche} d'autant plus qu'elle est prise sur un plus grand cercle.
\note{« circulaire » est écrit au-dessus de la ligne, sans signe d'insertion net.}

Proposition

Expliquer pourquoy se rompent les rayons par la rencontre d'un diaphane different de celuy d'où ils \uncertain{partent}, ou par \struck{\ill{}}
\note{énoncé souligné d'un trait continu, la fin biffée.}

Puisque nul ne doute que la lumiere ne se rompe lors qu'elle rencontre un milieu different en transparence, \& d'ailleurs \struck{\ill{}} \uncertain{c'est} dans l'onziesme prop. de l'optique suivant l'experience qui sert de fondement à nostre discours, \add{à sçavoir} \struck{\ill{}} \uncertain{nous} n'expliquerons la cause, comme \struck{\ill{}} avons fait en parlant de la reflexion : \& puis nous \uncertain{dirons} \struck{quelle est} \uncertain{precisement} les loix de cette refraction, c'est àdire à quels angles elle se fait \struck{suivant la \ill{}} dans les differentes inclinations du rayon entrant dans le diaphane, ou en sortant des mesmes diaphanes, car il se rompt \struck{\ill{}} aussi bien en sortant qu'en entrant.
\note{« l'onziesme prop. de l'optique » : lecture offerte avec doute ; la ligne est une addition interlinéaire serrée. Le texte s'arrête ici ; le bas de la page est blanc.}

\page{38}
\note{Ouverture : à gauche, une feuille de papier blanc sans écriture ; au milieu, deux onglets de montage ; à droite, le recto d'un feuillet de la main des vues 21–36, écrit sur son tiers supérieur seulement. En haut à droite, « 30 », rond et appuyé, et plus haut « 39 », d'un trait plus fin. Le cachet de la bibliothèque est apposé dans le texte. La page de gauche de la vue 37, verso du feuillet paginé 28, et la page de droite de la même vue, blanche et numérotée 29, ne portent que la transparence du recto.}

Proposition XXII.

Expliquer pourquoy si l'air fait faire refraction à la lumiere, \& si elle peut passer par le vuide.
\note{énoncé souligné d'un trait continu.}

La premiere partie de cette proposition appartient à la dioptrique, c'est pourquoy je n'en parle icy qu'en general, \& dis seulement qu'elle se rompt en y tombant obliquement, \uncertain{si l'esprit} \uncertain{premier} \uncertain{certain} par où passe la lumiere du soleil, ou des autres astres, par exemple l'ether qui est depuis le soleil jusques à nostre atmosphere, est plus subtil \& moins opaque \struck{\ill{}} ou plus diaphane que l'air ; de sorte qu'il se peut faire que le rayon se rompe plusieurs fois depuis sa source, jusques à ce qu'il arrive à nous suivant les differentes vapeurs ou atmospheres qu'il peut rencontrer, \& specialement à ceux des astres qui sont plus proches de nous que le soleil.

Mais parceque nous ne pouvons parler de ce qui est ailleurs, \& que quand nostre atmosphere seroit \uncertain{cause} de la moitié des rayons solaires \struck{\ill{}} \uncertain{refraction} dans les luminaires, \uncertain{nous ne} \uncertain{pourrions} \uncertain{la} trouver, je laisse cette partie à la \uncertain{subtile} \uncertain{curiosité} : si ce n'est \uncertain{qu'on} \uncertain{veuille} \uncertain{par} \uncertain{ce} \add{\ill{}} \uncertain{calculer} la force du soleil du midy \struck{\ill{}} quand le soleil est vertical, comparée à celle du matin ou du soir, \& des autres heures du jour, \add{quand} le soleil nous frappe obliquement.
\note{le passage « quand nostre atmosphere … la trouver » est écrit vite et chargé ; la lecture n'en donne que l'ossature.}

Quant à l'autre partie,
\note{le texte s'arrête là, au milieu de la page, sur une virgule ; le reste du feuillet est blanc. Rien de la suite n'est dans ce lot.}

\page{39}
\note{Page de droite : une feuille de garde moderne, sans texte ancien, qui sépare deux parties du volume. Au milieu, à l'encre, sur deux lignes, « Pascal / autographe », et au-dessous, d'une autre encre ou au crayon, « et copies ». En bas à gauche, « Libri 1848 » ; en bas à droite, « FR. nouv. acq. 5176 ». En haut à droite, un grand « 35 » tracé d'un trait large et, au-dessous, « 31 », rond et appuyé. La page de gauche de la vue est le verso blanc du feuillet de la vue 38, le recto y transparaissant. Ce feuillet est une marque de classement, non un texte : il est signalé ici parce qu'il annonce ce qui suit, et rien de ce qu'il dit n'est repris comme attribution.}

\page{40}
\note{Page de droite : un feuillet d'une autre main que tout ce qui précède dans ce lot, une cursive ronde, penchée, régulière, à l'encre noire, sur un papier bruni, taché et froissé ; le coin supérieur droit, sombre et plissé, cache la fin de plusieurs lignes. En haut à droite, « 32 », d'un trait appuyé, le 3 surélevé, et à sa gauche, plus bas, « 23 » ; voir l'en-tête. La page de gauche de la vue est une feuille blanche, sans écriture. Le cachet de la bibliothèque (« R.F. ») est apposé dans la marge de gauche. Dans la même marge, à mi-hauteur, une petite bande de papier collée porte une note de neuf ou dix lignes, très serrée, qui n'a pas été lue. Le sujet est la règle des partis au jeu. Rien n'y nomme l'auteur.}

\ill{} \uncertain{Sy} \ill{} Jeux a \uncertain{Brin} \uncertain{J'ay} \uncertain{les} $\frac{3}{4}$ de tout l'argent. partant Il y a a \ill{} 3 contre \uncertain{1} que \ill{}
\note{ligne ajoutée en tête de la page, au-dessus du titre, d'une encre plus noire, dans la partie brunie et plissée du bord ; plusieurs mots sont perdus dans la tache. Sous son début, un mot seul, « Espargne » ou « Esperance », n'a pas été lu sûrement.}

Regle \struck{\ill{}} \add{auxquelles} se peuvent rapporter les partis
\note{titre, centré. « auxquelles » est écrit au-dessus d'un mot biffé.}

Sy mon \struck{\ill{}} \add{hasard} est egale d'y gagner (par Ex.) 10, ou 12 --- cela me vaut 11 --- c'est à dire tousjours la moitié des deux \struck{\ill{}} \add{hasards} ou esperances.
\note{« hasard » et « hasards » sont écrits au-dessus de mots biffés. Le chiffre après « vaut » paraît « 11 » par le sens ; il est écrit en deux traits courbes et peut se lire autrement.}

Et ainsy quand on propose. vous aurés 10, ou vous aurés 12, Il m'appartiendra 11. ce qui \uncertain{fait} \ill{} \ill{}
qui est que mon esperance dans le party vaut une telle somme ou nombre, qu'y ayant \uncertain{un} \uncertain{semblable} \ill{}
et le jouant contre \uncertain{luy} egal, \struck{\ill{}} \add{a condition} que le gagnant donne au perdant la plus petite somme, je me \uncertain{tiendray} \ill{}
tousjours dans la mesme esperance d'avoir le mesme party qu'on m'a proposé.
\note{les fins de ces trois lignes disparaissent dans le coin bruni et plissé du feuillet. « a condition » est écrit au-dessus d'un mot noirci.}

Par Ex. on me donne a choisir au hasard entre deux sommes cachées, l'une de 9, l'autre de 13. \ill{}
cela vaut autant que 11. parceque Jouant 11 contre 11 a condition que le gagnant donne 9 au perdant
je reviens a la mesme esperance que j'avois auparavant qui est d'avoir tousjours 9 ou 13. car sy je
gagne j'auray 22 dont donnant \struck{\ill{}} 9 au perdant j'auray \struck{\ill{}} encor 13. sy je perds j'ay les 9 que
me \uncertain{revient}.
\note{« 22 dont donnant 9 » : le chiffre avant « au perdant » est biffé et récrit.}

Mais plus court, et moins embrouillé. Il n'y a qu'a dire que des deux partis proposés
j'ay desja seurement le moindre, et Il n'y a qu'a partager par moitié l'excedant du plus grand \ill{}
moindre, et l'adjouster au moindre --- pour faire a \uncertain{chacun} ce qui \uncertain{leur} revient.
\note{« l'excedant du plus grand » : la fin de la ligne est prise dans le pli du coin droit.}

ainsy quand j'ay hasard egal a tirer 9 ou 13. j'ay desja seurement 9 quoy qu'arrive, Et
le hasard est egal d'avoir aussy tost 13 que 9, sy on veut quitter sans que je tire Et partagions a\ill{}
est deu a chacun, Il est certain que pouvant avoir aussy tost 13 que 9 Il faut partager entre nous les \ill{}
dont 13 excede 9, qui seront 2 pour moy, Et avec 9 font 11. ce qui est le mesme que l'autre Regle.
\note{les derniers mots des lignes sont pris dans le coin bruni ; le mot perdu à la fin de l'avant-dernière ligne est le nombre dont 13 excède 9.}

de mesme si le hasard n'est pas egal et que j'aye par exemple 2 pour avoir 13 contre
1 pour n'avoir que 9 il faut partager l'excedent en 3 parties et m'en donner deux \&c.
\note{ces deux lignes sont d'une autre encre et d'un trait plus gros, plus droit, sans les liaisons de la main du feuillet : une addition, peut-être d'une autre main, serrée entre le paragraphe et le titre suivant.}

Exemples suivant ces Regles

En 2 Jeux combien vaut le 1\textsuperscript{er}.
\note{titre et sous-titre, centrés.}

Je joue contre quidam 6 escus en 2 jeux. J'ay le premier jeu, Et nous voulons quitter, combien me \ill{}
il. c'est 9. car sy je faisois encor \struck{\ill{}} \add{le} jeu suivant j'ay 12 Et luy rien --- s'il le fait j'ay encor 6, Et luy \ill{}
or le hasard est egal a chacun pour ce jeu la. donc mon esperance est egale pour avoir 12, ou pour
avoir 6. --- donc partageant les deux esperances par moitié j'ay 9, Et luy 3. (qui est la premiere Regle) ou \ill{}
j'ay desja seurement 6, Et j'ay hasard egal pour les 6 autres --- donc en quittant le jeu Il me revient la \ill{}
de les 6 autres qui est 3 Et joint a 6 font 9 (qui est la 2\textsuperscript{e} Regle.) --- Et pour achever de rapporter \ill{}
chose a la premiere Regle Il faut dire, sy je joue 9 contre 9 a condition que celuy qui gagnera
donnera 6 a l'autre, Il est certain que je reviens encor a ma premiere esperance d'avoir 12, ou 6.
\note{les fins de lignes sont prises dans le bord bruni du feuillet. « 1\textsuperscript{er} » : « le 1\textsuperscript{er} » jeu. Les parenthèses sont de l'auteur.}

or puisque j'ay 9 contre 3 sy un autre veut prendre mon jeu Il me doit donner 9 --- ou sy \uncertain{voulant}
\uncertain{vous} gager que je gagneray tout, Il doit mettre 3 contre 9.

Pour d'autres exemples a plusieurs Jeux Il faut prendre premierement les Regles suivantes

I Regle. a chaque jeu chacun joue jeu egal, dans un hasard egal.

2 \struck{Celuy \ill{}} \struck{Pour} voir combien revient a chacun quand on joue en plusieurs jeux
Il ne faut que reduire la partie en sorte qu'il faille 2 jeux a l'un pour
gagner, et a l'autre seulement un, auquel cas celuy cy a les $\frac{3}{4}$ de la
somme. Et suivre apres cela les Regles precedentes.
\note{« Celuy \ill{} » est biffé en tête de la règle 2, et le mot suivant l'est aussi. La fraction $\frac{3}{4}$ est à la marge, prise dans le pli.}

Exemple --- on joue en 3 jeux 8 escus contre 8 autres

Le premier jeu vaudra $\frac{3}{16}$ si A le gagne Il luy revient $\frac{11}{16}$ --- Et a B seulement $\frac{5}{16}$

Le second jeu vaut encor $\frac{3}{16}$ --- A le gagnant a $\frac{14}{16}$ Et B $\frac{2}{16}$

Le 3. jeu vaut $\frac{2}{16}$
\note{les dénominateurs ont un premier chiffre à trait d'attaque qui se lit aussi bien 2 que 1 ; les sommes des deux parts ($11 + 5$, $14 + 2$, $12 + 4$) donnent 16, et c'est la lecture retenue.}

Sy A le gagne Il a tout

Sy B. A a encor $\frac{12}{16}$ Et B $\frac{4}{16}$

La mesme par les Regles est cecy
\ill{} A a 2 \ill{} a \ill{} on veut sçavoir ce qui luy revient \ill{} Il faut dire \ill{}
\note{les deux dernières lignes courent sur le bord inférieur du feuillet, bruni et plissé ; seuls quelques mots se lisent. La suite, s'il y en a une, est au verso, hors de ce lot (vue 41).}

\end{document}
